\section{Related Work}
\label{sec:related_work}

Our work is closely related to three lines of research:
(i) spectral anisotropy in language model representations,
(ii) low-bit quantization under anisotropic spectra,
and (iii) optimizer design under structured gradient geometry.
We review each in turn and clarify how our contribution differs.

\paragraph{Anisotropy and outlier dimensions in language models.}
A growing body of work has documented that contextual representations produced
by Transformer language models are highly anisotropic.
Early studies showed that word and sentence embeddings concentrate in narrow
cones under cosine similarity, reducing representational diversity
\citep{arora2017simple,mu2017allbutthetop,ethayarajh2019contextual}.
Subsequent analyses identified small sets of ``rogue'' or ``outlier'' dimensions
that dominate similarity computations and downstream behavior
\citep{timkey2021rogue,kovaleva2021bertbusters,puccetti2022frequency}.
More recent work argues that anisotropy is closely tied to frequency skew,
syntactic structure, and training dynamics rather than being an inherent flaw
of the Transformer architecture \citep{rudman2023stable,machina2024anisotropy}.

Our work builds on these observations but shifts the focus from representation
quality at inference time to \emph{training-time spectral geometry}.
Rather than treating anisotropy as a nuisance to be removed, we show that the
coherent structure underlying anisotropy can be exploited algorithmically.

\paragraph{Low-bit quantization and spectral decomposition.}
Quantization of large language models has progressed rapidly from FP32 to BF16,
FP8, and recently FP4, motivated by substantial gains in efficiency
\citep{micikevicius2022fp8,peng2023fp8,perez2023fp8}.
However, extreme low-bit regimes exacerbate sensitivity to wide numerical
distributions and outliers.
A range of methods address this problem via reparameterization or preprocessing,
including channel-wise scaling and shifting \citep{xiao2023smoothquant,wei2023outlier},
orthogonal or Hadamard rotations \citep{suresh2017distributed,ashkboos2024quarot,ashkboos2025halo},
and outlier separation into high-precision branches \citep{li2024svdquant}.

Metis~\citep{cao2025metis} identifies spectral anisotropy as the \emph{fundamental}
barrier to FP4 training and proposes explicit spectral decomposition of weights,
activations, and gradients to separate dominant singular components from the
long tail.
By quantizing these components independently, Metis enables robust W4A4G4
training with minimal accuracy loss.
However, Metis still relies on estimating a low-rank spectral subspace, using
random sampling and projection to amortize the cost of SVD.

Our work is complementary but conceptually distinct.
Rather than approximating the spike subspace via randomized spectral methods,
we show that, due to the statistical structure of language data, a shared
linguistic component can be isolated by a much simpler operation—batch
averaging—without explicit decomposition.
This provides a lightweight surrogate that can fill the role of SVD-based spike
suppression in quantization-sensitive settings.

\paragraph{Optimizers under anisotropic gradient geometry.}
Adaptive optimizers such as Adam and AdamW \citep{kingma2014adam,loshchilov2017adamw}
operate coordinate-wise and ignore correlations within matrix-valued parameters.
When gradient energy concentrates in a low-rank spike, their second-moment
statistics become dominated by these directions, suppressing updates in the
long-tail subspace.
Matrix-aware optimizers, including K-FAC \citep{martens2015kfac}, Shampoo
\citep{gupta2018shampoo}, and SOAP \citep{vyas2024soap}, capture structured
correlations but incur significant computational and memory overhead.

Spectra~\citep{spectra2026} directly targets this issue by identifying and
suppressing the dominant low-rank spike subspace in gradient updates, while
leaving the noise-sensitive tail unamplified.
Spectra tracks the spike using cached power iteration and applies localized
spectral shaping with low overhead.

Our contribution is orthogonal to optimizer design.
We provide a data-geometric explanation for why the dominant spike is
\emph{coherent and shared across tokens}, and show that averaging removes a
direction lying almost entirely within the spike subspace.
This insight complements spike-aware optimizers such as Spectra by providing a
simple, hardware-friendly primitive that attenuates spike-dominated statistics
without requiring iterative spectral routines.

\paragraph{Efficient low-rank approximation.}
Classical SVD and QR-based methods are expensive in large-scale training loops.
Randomized and streaming techniques offer scalable alternatives for estimating
dominant subspaces \citep{halko2011randomized,tropp2012user,oja1982pca}.
Both Metis and Spectra exploit these tools to reduce the cost of spectral
interventions.

Our work shows that, under the specific statistical structure of human language
corpora, even these approximations may be unnecessary for isolating the most
harmful shared component.
By leveraging coherence in linguistic structure and incoherence in semantic
residuals, averaging acts as a conservative but effective spectral
approximation.